\documentclass[11pt,a4paper]{article}
\usepackage[margin=1in]{geometry}
\usepackage{hyperref}
\usepackage{enumitem}
\usepackage{graphicx}
\usepackage{xcolor}

% -----------------------------------------------------
% bvraghav-tiet-ucs503p-article.sty
% -----------------------------------------------------

% Variable: Lab Instructor
\makeatletter \newcommand{\BvrSetLabInstructor}[1]{%
  \gdef\Bvr@LabInstructor{#1}}
% default
\newcommand{\Bvr@LabInstructor}{Nisha Thakur}
% public accessor
\newcommand{\BvrLabInstructorValue}{\Bvr@LabInstructor}
\makeatother

% Add subtitle
\usepackage{titling}
% -- Set title bold
\pretitle{\begin{center}\bfseries\fontsize{18bp}{18bp}\selectfont}
\makeatletter
\newcommand{\BvrSubtitle}[1]{%
  \posttitle{%
    \par\end{center}
    \begin{center}\large#1\end{center}
    \vskip 0.5em}%
}
\makeatother

% Hints in the section title(s)
\newcommand{\BvrHint}[1]{%
  {\normalfont\normalsize\itshape\hfill #1}}

% -----------------------------------------------------

\title{Smart Photo Drive}

\BvrSetLabInstructor{Ms. Nisha Thakur}

\BvrSubtitle{%
  Project Proposal for UCS503 -- Software Engineering  \\
  Submitted to: \BvrLabInstructorValue \\ \bfseries
  Thapar Institute of Engineering and Technology%
}

\author{
  Harshjyot Kaur -- 1024031176
  \and
  Jaiveer Singh -- 1024030976
  \and
  Archita Sharma -- 1024030715
  \and
  Avneet Kaur -- 1024031173
}

\date{\today}

\begin{document}
\maketitle

\section{Higher-order goal%
  \BvrHint{\ldots secondary framing}}

This project aims to improve the way photographs captured during large
events are searched, organized, and accessed by individuals. Event
photographers commonly capture thousands of photographs, making it
difficult for a participant to manually locate all photographs in which
they appear. The higher-order goal of Smart Photo Drive is to transform
this time-consuming manual search into an automated and user-friendly
photo retrieval process.

The immediate engineering objective is to develop a deployable web-based
system that uses artificial intelligence for face-based photo retrieval.
A user should be able to provide a reference photograph containing their
face, after which the system searches the event photograph collection
and presents the photographs that are likely to contain the same person.

\section{Time-to-value %
  \BvrHint{\ldots primary heuristic}}

\begin{itemize}[leftmargin=0.7cm]
    \item We will first implement the core photo retrieval workflow:
    event photo upload, reference face input, face processing, matching,
    and display of retrieved photographs.

    \item The project will be developed incrementally so that a working
    prototype can be demonstrated early, followed by improvements in
    matching quality, usability, performance, and deployment.

    \item We will prioritize a functional end-to-end system over
    unnecessary complexity, allowing the team to validate the core idea
    before adding advanced features.

    \item Success in the initial iteration is defined by the system's
    ability to process an event photo collection and retrieve relevant
    photographs for a given user.
\end{itemize}

\section{Problem Statement}

Event photographers may capture hundreds or thousands of photographs
during college events, weddings, conferences, festivals, sports events,
and other large gatherings. After the event, participants often need to
manually browse through large collections of photographs to find the
images in which they appear.

Common pain points include:

\begin{itemize}[leftmargin=0.7cm]
    \item \textbf{Manual searching:} Users may have to browse through
    hundreds or thousands of photographs individually to locate their
    pictures.

    \item \textbf{Time consumption:} Searching large event collections
    manually can require significant time, especially when the user
    appears in many photographs.

    \item \textbf{Poor discoverability:} Conventional photo folders
    generally organize images by filename, timestamp, or directory
    rather than by the people appearing in the photographs.

    \item \textbf{Large-scale collections:} As the number of photographs
    increases, manually searching the collection becomes increasingly
    inconvenient.

    \item \textbf{Repeated effort:} Multiple users searching for their
    own photographs may repeatedly perform the same manual browsing
    process.
\end{itemize}

\textbf{Why this matters now (contextual relevance):} Modern events
generate increasingly large digital photo collections. An automated
face-based retrieval system can reduce the effort required to locate
personal photographs while providing photographers or event organizers
with a more efficient way to distribute photographs.

\section{Proposed Solution %
\BvrHint{\ldots Software Proposition}}

\subsection{Overview}

Build a \textbf{web-based AI-powered photo retrieval system} that allows
event photographers or administrators to upload a collection of event
photographs and allows users to retrieve photographs containing their
face.

The proposed system will contain the following components:

\begin{itemize}[leftmargin=0.7cm]
    \item \textbf{Event photo management:} an administrator or
    photographer uploads a large collection of photographs associated
    with an event.

    \item \textbf{Reference photo input:} a user provides a photograph
    containing their face to serve as the reference for the search.

    \item \textbf{Face detection and processing:} the system identifies
    faces in the reference image and event photographs using an
    appropriate computer vision model.

    \item \textbf{Face matching:} facial representations are compared
    to identify photographs that are likely to contain the same person.

    \item \textbf{Photo retrieval:} matching photographs are collected
    and presented to the user through the web interface.

    \item \textbf{Photo access:} users can view their matched
    photographs and download or access them as permitted by the
    application.
\end{itemize}

The project will evaluate suitable AI-based face detection and
recognition techniques during implementation. YOLO-based face detection
may be considered as one candidate approach, while the final model and
matching method will be selected based on accuracy, processing time,
availability of suitable models, and deployment feasibility.

\subsection{Core workflow}

\begin{enumerate}[leftmargin=0.7cm]
    \item The photographer or administrator uploads the photographs
    captured during an event.

    \item The system processes the uploaded photographs and detects
    faces present in the images.

    \item A user opens the web application and provides a reference
    photograph containing their face.

    \item The system processes the reference face and compares it with
    the facial information extracted from the event photographs.

    \item Photographs with sufficiently similar facial representations
    are identified as candidate matches.

    \item The system displays the retrieved photographs to the user,
    who can view and download the relevant images.
\end{enumerate}

\subsection{Operational constraints for an educational deployment}

\begin{itemize}[leftmargin=0.7cm]
    \item The system should provide a responsive web interface that can
    be accessed from commonly available computers and mobile devices.

    \item The solution should use established computer vision and web
    development technologies instead of requiring the development of a
    novel face-recognition algorithm.

    \item The initial implementation should be designed for a realistic
    educational project scale while remaining extensible to larger
    event photo collections.

    \item User photographs and event photographs should be handled with
    appropriate access controls and privacy considerations.
\end{itemize}

\section{Solution Approach %
  \BvrHint{\ldots Engineering focus}}

This is an engineering course project, so the primary focus will be on
developing a complete, reliable, and deployable software system rather
than proposing a new face-recognition research algorithm.

The proposed implementation will use a modern web architecture
consisting of a ReactJS frontend, a Python-based FastAPI backend, and
MongoDB for application data. The AI processing component will be
integrated with the backend and will use an appropriate face detection
and recognition pipeline selected during development.

\subsection{Face detection and recognition %
  \BvrHint{\ldots AI component}}

The AI component will perform two major tasks:

\begin{itemize}[leftmargin=0.7cm]
    \item \textbf{Face detection:} identify the location of faces within
    the reference image and event photographs.

    \item \textbf{Face matching:} generate or obtain suitable facial
    representations and compare the reference face with faces found in
    the event photographs.
\end{itemize}

The team will evaluate practical face detection and recognition
approaches based on accuracy, computational requirements, inference
speed, and ease of integration with the web application. YOLO-based
models may be evaluated for the face detection component, while a
suitable face embedding or recognition method may be used for matching.

The system will use a similarity threshold to determine candidate
matches. The threshold will be tuned using test photographs to balance
the number of correctly retrieved photographs against incorrect
matches.

\subsection{Web architecture %
  \BvrHint{\ldots web-based solution}}

A typical three-tier architecture will be used:

\begin{itemize}[leftmargin=0.7cm]
    \item \textbf{Frontend:} ReactJS will provide the user interface for
    event management, reference photo upload, search initiation, and
    display of retrieved photographs.

    \item \textbf{Backend API:} FastAPI will provide APIs for image
    upload, processing requests, face-matching operations, retrieval of
    results, and application logic.

    \item \textbf{Database:} MongoDB will store application metadata
    such as users, events, uploaded photo information, processing
    status, and references required for retrieving photographs.
\end{itemize}

The actual image files may be stored using suitable file or object
storage, while the database will maintain the corresponding metadata
and references.

\subsection{AI processing pipeline}

The proposed processing pipeline is:

\begin{enumerate}[leftmargin=0.7cm]
    \item Upload the event photographs.
    \item Validate and preprocess the images.
    \item Detect faces in the event photographs.
    \item Generate facial representations for detected faces.
    \item Store the required processing information for later matching.
    \item Accept a user's reference photograph.
    \item Detect and represent the user's face.
    \item Compare the reference representation against stored
    representations.
    \item Retrieve photographs containing sufficiently similar faces.
\end{enumerate}

This approach avoids repeatedly performing the complete image-processing
pipeline whenever a new user searches the same event collection.

\subsection{Development and deployment approach}

The project will be developed incrementally using version control and
regular testing. The team will maintain the project source code in a
GitHub repository and use separate development tasks for frontend,
backend, database, AI processing, testing, and deployment.

Where practical, automated build and testing steps will be introduced
during development. The final system will be deployed to a suitable
web-accessible environment so that the complete photo retrieval
workflow can be demonstrated.

\section{Evaluation Criterion %
  \BvrHint{\ldots measurable and attributable}}

The project will evaluate the system using measurable criteria related
to photo retrieval quality, processing performance, and usability.

\subsection{Primary evaluation metric}

\textbf{Photo retrieval precision:} the proportion of photographs
returned by the system that actually contain the queried user.

\begin{itemize}[leftmargin=0.7cm]
    \item The system will be tested using a labelled set of event
    photographs where the actual presence of the target person is known.

    \item The retrieved results will be compared with the expected
    photographs to measure the quality of the matching process.

    \item The similarity threshold will be adjusted using validation
    photographs to obtain a practical balance between correct and
    incorrect matches.
\end{itemize}

\subsection{Secondary evaluation metrics}

\begin{itemize}[leftmargin=0.7cm]
    \item \textbf{Recall:} percentage of the target user's actual
    photographs that are successfully retrieved.

    \item \textbf{Retrieval time:} time taken by the system to return
    matching photographs after a user's search request.

    \item \textbf{Processing throughput:} number of event photographs
    that can be processed within a given period.

    \item \textbf{False-match rate:} proportion of retrieved
    photographs that do not contain the target user.

    \item \textbf{User experience:} ease with which a user can upload a
    reference image, search the collection, and access the retrieved
    photographs.

    \item \textbf{System reliability:} successful completion of photo
    upload, processing, search, and retrieval operations during
    testing.
\end{itemize}

\subsection{Pilot validation plan %
  \BvrHint{\ldots high-level}}

\begin{itemize}[leftmargin=0.7cm]
    \item Prepare a representative event photo collection containing
    multiple people and different photographic conditions.

    \item Select a small group of test users and provide each user with
    one or more reference photographs.

    \item Compare system-generated results against manually verified
    photographs to evaluate precision and recall.

    \item Measure retrieval time for different collection sizes to
    understand how system performance changes as the number of
    photographs increases.

    \item Collect basic user feedback regarding the simplicity and
    usefulness of the retrieval workflow.
\end{itemize}

\section{Scalability %
  \BvrHint{\ldots with foresight}}

The proposed system should be capable of handling growth in both the
number of event photographs and the number of users.

\begin{enumerate}[leftmargin=0.7cm]

    \item \textbf{Scalable (theoretically):} The system can support
    larger photo collections by:
    \begin{itemize}[leftmargin=0.7cm]
        \item separating frontend, backend, and AI processing
        responsibilities,
        \item storing photo metadata separately from image files,
        \item indexing frequently accessed database fields,
        \item processing large photo collections in batches,
        \item avoiding repeated face detection for photographs that
        have already been processed.
    \end{itemize}

    \item \textbf{Operationally deployable (practically):} The system
    should be deployable using commonly available hosting and managed
    services:
    \begin{itemize}[leftmargin=0.7cm]
        \item ReactJS frontend deployment,
        \item FastAPI backend deployment,
        \item managed MongoDB database,
        \item suitable image/object storage,
        \item containerized or platform-hosted deployment where
        appropriate.
    \end{itemize}

\end{enumerate}

\section{Engine availability heuristic}

A mature set of tools and services is available to implement the
proposed Smart Photo Drive system:

\begin{itemize}[leftmargin=0.7cm]
    \item \textbf{ReactJS} for the responsive web frontend.

    \item \textbf{FastAPI} for backend REST APIs and integration of the
    AI processing pipeline.

    \item \textbf{MongoDB} for storing users, event information,
    photograph metadata, and processing-related information.

    \item \textbf{Computer vision and face-recognition models} for face
    detection and matching.

    \item \textbf{Python-based AI ecosystem} for image preprocessing,
    model inference, and facial representation generation.

    \item \textbf{Git and GitHub} for version control and collaborative
    software development.

    \item \textbf{Cloud or platform-based deployment services} for
    making the application accessible through the web.
\end{itemize}

These mature technologies allow the team to focus primarily on system
integration, workflow design, photo retrieval accuracy, usability,
testing, and deployment rather than developing all underlying
infrastructure from scratch.

\section{Project scope and deliverables %
  \BvrHint{\ldots iteration-friendly}}

\subsection{Initial deliverable %
  \BvrHint{\ldots first iteration}}

The first implementation iteration will focus on establishing the
complete basic retrieval pipeline:

\begin{itemize}[leftmargin=0.7cm]
    \item Basic ReactJS web interface.
    \item Event photo upload functionality for the administrator or
    photographer.
    \item Reference photograph upload for users.
    \item Basic image validation and preprocessing.
    \item Face detection in uploaded photographs.
    \item Initial face representation and matching pipeline.
    \item Retrieval and display of matching photographs.
    \item FastAPI backend APIs connecting the frontend with the
    processing pipeline.
    \item MongoDB integration for storing application metadata.
\end{itemize}

\subsection{Subsequent deliverables}

\begin{itemize}[leftmargin=0.7cm]
    \item Improved face matching accuracy through model and threshold
    evaluation.
    \item Support for larger event photo collections.
    \item User and administrator authentication where required.
    \item Improved photo search and filtering.
    \item Photo download and sharing functionality.
    \item Processing-status indicators for large photo uploads.
    \item Performance optimization for repeated searches.
    \item Automated testing of frontend, backend, and AI-processing
    components.
    \item Deployment of the complete system to a web-accessible
    environment.
    \item Documentation of the architecture, APIs, testing, and
    deployment process.
\end{itemize}

\section{Risks and mitigations}

\begin{itemize}[leftmargin=0.7cm]

    \item \textbf{Face recognition accuracy:} Differences in lighting,
    pose, facial expressions, image quality, occlusion, and camera
    angles may produce incorrect or missed matches. The system will be
    evaluated using representative test images and the matching
    threshold will be tuned accordingly.

    \item \textbf{Large photo collections:} Processing thousands of
    photographs may require substantial computation and time. Face
    information can be precomputed during event-photo processing so
    that subsequent user searches do not require complete
    reprocessing.

    \item \textbf{Privacy and consent:} Facial information and
    photographs are sensitive user data. The system will minimize
    unnecessary data collection, restrict access to photographs, and
    provide appropriate controls for uploaded data.

    \item \textbf{Storage requirements:} Large collections of
    photographs may require significant storage. Image files and
    metadata will be managed separately so that the database does not
    unnecessarily store large binary files.

    \item \textbf{Model and deployment limitations:} AI inference may
    require more computational resources than ordinary web requests.
    The team will evaluate model size and inference requirements while
    selecting an approach suitable for the educational deployment.

    \item \textbf{Evaluation ambiguity:} Retrieval quality can be
    difficult to judge without ground-truth data. A labelled test
    collection and manually verified expected results will be used for
    precision and recall evaluation.

    \item \textbf{Integration risk:} Frontend, backend, database, and
    AI components may introduce integration issues. The system will be
    developed incrementally with clearly defined APIs and regular
    integration testing.
\end{itemize}

\section{Summary}

This project proposes \textbf{Smart Photo Drive}, a deployable web-based platform designed to
automatically find photographs of a user from large event photo
collections.

The system will allow a photographer or administrator to upload event
photographs and allow users to provide a reference photograph
containing their face. An AI-based face detection and recognition
pipeline will then process the collection, identify candidate
photographs containing the user, and present the results through a
web-based interface.

The project meets the software engineering objectives by emphasizing
incremental development, measurable evaluation criteria, modular
architecture, scalability, testing, version control, and deployment.
The proposed ReactJS frontend, FastAPI backend, MongoDB database, and
computer vision pipeline provide a practical foundation for developing
the system while allowing the specific face-recognition model to be
selected and evaluated during implementation.

The final objective is to replace the traditional manual process of
searching through thousands of event photographs with a faster,
automated, and user-friendly photo retrieval experience.

\end{document}
